\documentclass[11pt,a4paper]{article}
\usepackage{../shared/hanzocover}
\usepackage[utf8]{inputenc}
\usepackage[T1]{fontenc}
% Shared preamble for the compute-market series.
% One style, one vocabulary, inputted by every paper so the series reads as one voice.
\usepackage[margin=1in]{geometry}
\usepackage{booktabs}
\usepackage{amsmath}
\usepackage{amssymb}
\usepackage{amsthm}
\usepackage{graphicx}
\usepackage{xcolor}
\usepackage{hyperref}
\hypersetup{colorlinks=true, linkcolor=blue, urlcolor=blue, citecolor=blue}
\usepackage{enumitem}
\usepackage{listings}
\usepackage{array}
\usepackage{tabularx}
\usepackage{siunitx}
\sisetup{group-separator={,}, group-minimum-digits=4}

\definecolor{kw}{rgb}{0.13,0.29,0.53}
\definecolor{cmt}{rgb}{0.40,0.40,0.40}
\definecolor{str}{rgb}{0.16,0.50,0.16}
\lstset{
  basicstyle=\ttfamily\footnotesize,
  keywordstyle=\color{kw}\bfseries,
  commentstyle=\color{cmt}\itshape,
  stringstyle=\color{str},
  showstringspaces=false,
  breaklines=true,
  columns=fullflexible,
  frame=single,
  framesep=4pt,
}

\newtheorem{definition}{Definition}
\newtheorem{proposition}{Proposition}
\newtheorem{remark}{Remark}

\newcommand{\code}[1]{\texttt{#1}}
\newcommand{\repo}[1]{\texttt{#1}}

% Series vocabulary — used identically in every paper.
\newcommand{\job}{\ensuremath{J}}
\newcommand{\bundle}{\ensuremath{R}}
\newcommand{\cost}{\ensuremath{\mathcal{C}}}

\tolerance=800
\emergencystretch=3em


\title{Byte Equality:\\Deterministic Work as a Market Primitive}

\author{
Hanzo AI Research\thanks{Corresponding author: research@hanzo.ai} \\
\textit{Hanzo Industries Inc (Techstars '17)} \\
Los Angeles, California \\
\texttt{https://hanzo.ai}
}

\date{2026}

\begin{document}
\hanzocoverpage

\begin{abstract}
The hardest problem in a decentralized compute market is not matching supply to
demand. It is establishing that the supplier computed what the buyer paid for.
Most designs answer this with replication or with succinct proofs, and both are
expensive: replication multiplies the compute bill by the committee size, and
proving multiplies it by three to six orders of magnitude.

For a large and underexploited class of work, neither is necessary. When a
computation is bit-reproducible, correctness is decidable by comparison against
a canonical implementation at the cost of running the canonical implementation
once --- and for a market that publishes fixed test vectors, at the cost of a
table lookup. We describe the byte-equality contract we enforce across a
29-primitive cryptographic corpus and three GPU backends, argue that it is a
market primitive rather than an engineering practice, and show that the
crossover threshold it produces as a by-product is itself a quantity the market
needs. We then draw the determinism boundary honestly: which AI work can be made
bit-exact, at what throughput cost, and where it cannot.
\end{abstract}

\tableofcontents
\newpage

\section{Two regimes of verification}

Let a job be a function $f$ and an input $x$; the provider returns $\hat{y}$ and
claims $\hat{y} = f(x)$. Verification splits on a single property.

\begin{definition}[Determinism class]
A job is in class $\mathsf{D}$ if $f$ is bit-reproducible: any correct
implementation on any conforming device returns byte-identical output for
identical input. It is in class $\mathsf{S}$ if $f$ is well-defined only up to a
tolerance --- floating-point reductions whose order is unspecified, sampling with
unpinned entropy, kernels whose tiling changes the summation order. It is in
class $\mathsf{O}$ if the buyer cannot characterise $f$ at all.
\end{definition}

The regimes have radically different verification economics, and a market that
does not distinguish them overpays for the easy case and underprotects the hard
one.

\begin{table}[h]
\centering
\begin{tabular}{llll}
\toprule
Class & Verification & Cost & Adjudication \\
\midrule
$\mathsf{D}$ & recompute canonical / KAT lookup & $O(1)$ amortised & proof of misbehaviour \\
$\mathsf{S}$ & replication or succinct proof & $m\times$ or $10^{3}$--$10^{6}\times$ & statistical or cryptographic \\
$\mathsf{O}$ & reputation, attestation, stake & --- & social \\
\bottomrule
\end{tabular}
\caption{Verification economics by determinism class.}
\end{table}

The point of this paper is that class $\mathsf{D}$ is much larger than it looks,
and that the machinery to exploit it already exists and is open source.

\section{The contract}

We hold every shipping GPU kernel to byte equality against a first-party CPU
canonical. The contract has three parts, and each exists for a distinct reason.

\paragraph{Equality on randomized inputs.} Every release, each primitive is
exercised on $N = 1000$ randomized inputs across the CPU canonical and every GPU
backend --- Metal, CUDA and WGSL. All outputs must be byte-identical. Random
inputs rather than fixed vectors matter because fixed vectors are what an
implementation is tuned against; randomization is what catches the tiling bug
that only fires on an odd tail.

\paragraph{Test-oracle isolation.} The audited upstream implementation used as
an oracle is never linked into the shipped library. This keeps the thing being
tested and the thing testing it genuinely independent, so a shared bug cannot
cancel out. Cross-oracle checking extends this: BN254 results are checked
against \code{gnark-crypto} (36 of 36 vectors) and independently against
\code{arkworks}, which are unrelated codebases and therefore give adversarial
distance rather than a second opinion from a cousin.

\paragraph{Memory hygiene.} Secret state does not leave attested accelerator
memory on the hot path. This is orthogonal to correctness but belongs in the
same contract because a market will run these kernels on machines the key owner
does not control.

\section{Why this is a market primitive}

An engineering team reads the above as a quality-assurance regime. A market
should read it as three distinct economic objects.

\subsection{Fraud becomes impossible rather than detectable}

Under replication, a dishonest provider is caught with probability depending on
committee composition, and adjudication is a vote. Under byte equality a
mismatch is a \emph{proof}: the canonical output is a published fact, and any
party can recompute the comparison. Slashing is then mechanical and requires no
quorum, no oracle, and no assumption about the fraction of honest providers.
This removes an entire class of governance from the protocol.

\subsection{Verification cost collapses}

Consider a market that would otherwise replicate three-fold for confidence. The
verification overhead is $200\%$ of the useful compute. Moving a job from class
$\mathsf{S}$ to class $\mathsf{D}$ removes that overhead entirely. On a market
where a nontrivial share of volume is cryptographic or integer-quantized work,
this is the single largest efficiency available, and it is available without any
new cryptography.

\subsection{The crossover threshold is a price signal}

A by-product of the same methodology is $N^{*}$, the batch size at which a GPU
kernel overtakes the single-thread CPU baseline for a given primitive. Measured
across the corpus, $N^{*}$ ranges from $1$ (homomorphic-encryption number
theoretic transform --- the GPU wins on a single call) to $1024$ (SHA-256 --- a
GPU is worse than a CPU until a thousand hashes are in flight).

That number is not a benchmark curiosity. It is the feasibility boundary of a
bid. A provider offering GPU capacity for a class-$\mathsf{D}$ job with batch
size below $N^{*}$ is offering something strictly worse than a commodity CPU at
a higher price. Publishing $N^{*}$ per primitive lets the market reject
mispriced supply structurally rather than discovering it through disappointed
buyers.

\begin{table}[h]
\centering
\begin{tabular}{lll}
\toprule
Primitive family & Measured speedup & Crossover $N^{*}$ \\
\midrule
Fused BLS aggregation      & $144.22\times$ & --- \\
FHE NTT (Metal, kernel-only, 10-iter median) & $16.92\times$ & $1$ \\
Batched Ed25519 verify (RFC 8032) at $N{=}4096$ & $26.7\times$ & --- \\
SHA-256                     & ---            & $1024$ \\
\bottomrule
\end{tabular}
\caption{Representative measurements from the corpus, single Apple M2 Ultra.
CPU NTT throughput at $N{=}2^{20}$ is 6--9 Mops/s on the same device.}
\end{table}

\section{The determinism boundary in AI work}

The natural objection is that neural inference is class $\mathsf{S}$ and
therefore none of this applies. That is true by default and false by
construction, and the distinction is worth money.

\subsection{What makes inference nondeterministic}

Floating-point addition is not associative. A reduction whose order depends on
block scheduling, an atomic accumulation whose interleaving depends on
occupancy, a flash-attention tiling chosen by an autotuner at load time --- each
makes output depend on hardware and scheduling rather than on $f$ alone. Batch
size changes reduction order, so even the same device gives different answers at
different loads. None of this is error; it is unspecified behaviour that was
never worth specifying.

\subsection{What makes it deterministic again}

Determinism is a kernel property one may choose to pay for:

\begin{itemize}[leftmargin=1.4em]
\item \textbf{Integer paths.} Weight-and-activation quantized inference in
      \code{int8} or \code{int4} with integer accumulation is exactly
      reproducible. There is no rounding order to disagree about.
\item \textbf{Batch-invariant kernels.} Fixing the reduction order independent of
      batch size and occupancy removes the load dependence, at a throughput cost
      typically in the low tens of percent.
\item \textbf{Pinned sampling.} Greedy decoding, or sampling from a specified PRF
      seeded by a value in the job specification, makes the token sequence a
      function of the request.
\item \textbf{Pinned numerics.} A declared accumulation dtype and a declared
      attention implementation, both carried in the job specification and
      attested by the provider.
\end{itemize}

The result is a job class we would write as $\mathsf{D}_{\text{inf}}$: slower
than unconstrained inference, and verifiable by recomputation on any conforming
device. Whether a buyer wants it is a price question, which is the right way for
it to be.

\subsection{Where it does not reach}

Training is not practically reproducible across heterogeneous hardware, and
large-batch \code{bf16} inference at maximum throughput is not either. Those
remain class $\mathsf{S}$ and need the tier ladder. The claim here is narrower
and defensible: a material fraction of paid compute --- all cryptographic work,
all quantized inference, all embedding generation with pinned numerics, all
search --- can be moved into class $\mathsf{D}$ deliberately.

\section{Proposal: determinism class as a first-class job attribute}

We suggest a market carry the class on the job rather than inferring it.

\begin{lstlisting}
job:
  kind:        inference | crypto | search
  determinism: D | D_inf | S | O
  numerics:    { dtype: int8, accum: int32, attention: fixed, sampling: greedy }
  canonical:   <content-addressed reference implementation>
  vectors:     <published KAT set, optional>
\end{lstlisting}

Three consequences follow immediately. Pricing differs by class because
verification cost differs by class, and the buyer chooses. Providers advertise
which classes they can honour, which is a capability constraint the allocator
already understands. And a class-$\mathsf{D}$ job needs no proof system at all,
which means a market can launch with real verifiable volume before any proving
infrastructure exists.

\section{What is available}

The corpus, the contract and the measurement harness are open source under
permissive licence: a 29-primitive library spanning hashes, AEAD, elliptic
curves including pairing-friendly ones, polynomial commitments, lattice
primitives, post-quantum signatures and key encapsulation, threshold schemes,
homomorphic-encryption building blocks and EVM precompile wrappers, with Metal,
CUDA and WGSL kernels held to the equality contract described above.

Adopting it does not require adopting anything else in our stack. The contract
is a testing discipline plus a published artifact; the artifact is what a market
needs, and it is reproducible by anyone who wants to check it.

\section*{Note on scope}

Nothing here establishes that a provider ran the job \emph{when} it said, on the
hardware it claimed, or without also serving someone else. Byte equality
establishes what was computed, not the circumstances. Those are separate
questions with separate mechanisms, treated in the companion papers on
verification tiers and on settlement.

\end{document}
